\chapter{STL 与 ranges}
\index{STL}\index{ranges}\index{view}\index{borrowed range}

\begin{introduction}
  \item 按访问模式、迭代器稳定性和所有权选择容器
  \item 用 ranges 算法和 view 组合行情处理步骤
  \item 解释 view 的惰性、借用关系与悬空风险
  \item 识别浮点归约对顺序的依赖
\end{introduction}

\chaptermeta{已完成第 4--6 章，会使用 vector、迭代器、lambda、类、引用和 RAII。}{Python 的 list comprehension、生成器和 dict 能表达相似管道，但 C++ 容器的存储、迭代器失效和 view 生命周期都必须显式推理。}{若管道结果异常，依次检查源容器是否仍存活、view 是否仍借用同一数据、算法是否改变了顺序，以及归约是否依赖浮点结合律。}

本章把第五章的行情分析器演化为可独立构建的成交管道。固定输入中 AAPL 有三笔成交，程序必须精确输出：

\begin{lstlisting}
symbol=AAPL trades=3 first_sell=102.00 low=100.00 high=102.00 gross=2520.00 net_quantity=15
\end{lstlisting}

完成判据还包括：非法买卖方向与缺失代码以状态 2 退出；惰性 view、悬空 view 和浮点归约顺序均有可运行证据；\path{code/ch07} 可从空构建目录独立配置。

\section{从访问模式选择容器}

\subsection{场景：时间序列与按代码查询不是同一种访问}

成交按到达顺序进入程序，适合由 \texttt{vector<Trade>} 连续保存；仓位却需要频繁按代码查询，若每次都从头扫描时间序列，意图和成本都不清楚。下面的程序同时保留时间序列，并用 \texttt{map<string, int>} 建立净数量索引：

\lstinputlisting[caption={按访问模式组合 vector、map 与算法}]{code/ch07/container_choices.cpp}

\texttt{vector} 拥有元素并提供连续存储、常数时间索引和高效尾部追加，是批量行情的默认候选；扩容仍会使指向旧元素的迭代器、指针和引用失效。\texttt{map} 也拥有键和值，按键排序并提供对数时间查找。它的下标运算在键缺失时先插入值初始化的整数，再执行复合赋值；只查询已有键时用 \texttt{at}，拼错键会抛出异常而不会悄悄插入。范围循环遍历 \texttt{map} 时按键顺序得到 AAPL、MSFT，因此输出末尾为 \texttt{order=AAPL,MSFT}。

只关心平均常数时间查找且无须键序时，可考虑 \texttt{unordered\_map}；频繁在两端插入时，可考虑 \texttt{deque}。容器选择取决于访问模式和失效规则。原始 \texttt{trades} 保持时间顺序，只有副本 \texttt{sorted} 交给 \texttt{ranges::sort}。这样“首笔卖出”和“最低价格”各自使用合适的顺序。

\texttt{ranges::find\_if(trades, predicate)} 直接接收范围，返回指向首个匹配元素的迭代器；没有匹配时等于 \texttt{trades.end()}。本例已知存在卖出才使用 \texttt{first\_sell->price}，项目版本会显式检查末端。

\begin{pythonbridge}
Python \texttt{list} 与 C++ \texttt{vector} 都保持顺序，\texttt{dict} 与关联容器都按键查询；不同点是 Python 引用通常不因列表扩容失效，而 C++ 迭代器直接关联具体存储。Python 的 \texttt{sorted} 返回新列表，C++ 的 \texttt{ranges::sort} 原地修改范围，因此需要先决定是否保留原顺序。
\end{pythonbridge}

失败诊断：若直接排序 \texttt{trades} 后再寻找“首笔卖出”，程序仍能编译，却把“首笔”偷换为“排序后的第一笔”。根因是一个容器同时承担时间线和排名两种语义；修复是保留原序列并排序副本。若 \texttt{find\_if} 可能失败，则在解引用前先与末端比较。

立即练习：把第一笔 AAPL 的代码改为 MSFT。预期 \texttt{aapl\_position=5}，最低价变为 \texttt{MSFT@100.00}，而行数、代码数、首笔 AAPL 卖出及 \texttt{order=AAPL,MSFT} 不变。

\section{view：惰性借用而非新容器}

\subsection{场景：先描述筛选，再决定何时遍历}

大量中间 \texttt{vector} 会复制元素并分配内存。C++20 view 可以先组合“保留哪些元素、如何转换”，真正迭代时才逐个求值：

\lstinputlisting[caption={filter 与 transform 的惰性管道}]{code/ch07/lazy_views.cpp}

管道运算符 \texttt{|} 把左侧范围交给右侧 adaptor。\texttt{filter} 保存谓词并跳过不满足条件的元素；\texttt{transform} 在解引用时计算新值。\texttt{selected} 只是一个 view，没有新建 \texttt{vector} 或复制价格。它借用 \texttt{prices}，所以 view 创建后把首价从 99 改为 110，稍后的遍历会看到 110，并输出 \texttt{selected=220.00,202.00}。

借用也意味着源范围必须覆盖 view 的所有使用。view 自身按值保存并不延长底层 \texttt{vector} 的生命；同时，若源容器发生使迭代器失效的修改，已有 view 的遍历同样可能失效。

\begin{pythonbridge}
Python 生成器表达式同样惰性，但生成器通常持有对源对象的 Python 引用，垃圾回收会延长对象生命。本例由左值 \texttt{vector} 建立的 C++ view 只借用源容器，不能替它保活数据；其他 view 是否拥有底层范围，取决于它包装的 range 或 view。两者都可能在迭代时看到源数据的后续修改。
\end{pythonbridge}

故意失败的 \path{code/ch07/failures/dangling_view.cpp} 从函数返回借用局部 \texttt{vector} 的 view。它不进入绿色目标；ASan 执行必须非零退出并报告 stack-use-after-return、stack-use-after-scope 或等价的栈访问类别。根因在于源范围过早结束生命，与 filter 本身无关。修复是让调用者拥有源容器、在源仍存活时消费 view，或把结果物化为拥有元素的容器。

立即练习：把过滤阈值从 100 改为 105。首价修改后仍通过，101 被排除，精确输出应变为 \texttt{selected=220.00}。

\section{把算法组合成行情管道}

\subsection{场景：过滤、排序、查找、转换与归约各守一个职责}

最终管道接收已解析的 \texttt{vector<Trade>} 与目标代码。核心实现如下，CSV 读取和 CLI 只负责边界：

\lstinputlisting[caption={成交筛选、排序、查找、转换与归约}]{code/ch07/src/pipeline.cpp}

第一步的 \texttt{selected} 借用参数 \texttt{trades}；函数执行期间参数引用有效。第二步用 \texttt{back\_inserter(sorted)} 把筛选结果追加到拥有元素的 \texttt{vector}，因为排序会修改元素顺序。空结果在调用 \texttt{front} 或 \texttt{back} 前被拒绝。

\texttt{find\_if} 在原时间线中找首笔卖出；\texttt{transform} 把每笔成交映射为 \(price \times quantity\)；\texttt{accumulate} 从 \texttt{0.0} 开始按迭代顺序累加，因此结果类型为 \texttt{double}。最后的 \texttt{map} 用正买入、负卖出累计每个代码的净数量。每个步骤的对象寿命都能从局部变量和参数边界读出。

\begin{pythonbridge}
Python 常把这段写成筛选列表、\texttt{sorted}、\texttt{next} 和 \texttt{sum}。C++ ranges 把“遍历什么”与“如何处理”分开，但不会自动把 view 物化，也不会自动检查末端。需要修改顺序或把结果带出源范围时，显式复制反而是清楚的所有权决定。
\end{pythonbridge}

失败诊断：把命令行代码改为 MSFT 时，过滤能得到一笔买入，但找不到卖出。程序以状态 2 输出 \texttt{error=no sell trade for symbol MSFT}，而不是解引用末端迭代器。查询 TSLA 则更早得到 \texttt{no trades for symbol TSLA}。错误发生在哪个前置条件，诊断就停在哪一层。

立即练习：保持数据不变，只把命令中的 AAPL 改为 MSFT。验证退出状态为 2、stdout 为空、stderr 精确包含 \texttt{error=no sell trade for symbol MSFT}。

\section{浮点归约有顺序}

\subsection{场景：可归约不等于满足结合律}

算法接口允许累加，不代表浮点加法等同于实数加法。下面只改变遍历顺序：

\lstinputlisting[caption={相同浮点值的两种归约顺序}]{code/ch07/reduction_order.cpp}

有限精度下，\(10^{16}+(-10^{16})\) 先抵消，再加 1 得到 1；反向遍历时，1 先与 \(-10^{16}\) 相加会被舍入，最终得到 0。因此精确输出为 \texttt{forward=1 reverse=0}。\texttt{accumulate} 的顺序确定，却仍依赖输入顺序；允许重排的并行归约可能产生另一条舍入路径。

\begin{pythonbridge}
Python 的普通 \texttt{float} 通常同样是 IEEE 754 双精度，\texttt{sum} 也不能把浮点数变回实数。NumPy 的实现可能使用不同分块或成对求和，因此结果与逐项 Python 循环不同并不自动说明其中一个“错了”；必须先定义容差与数值口径。
\end{pythonbridge}

失败诊断：这里没有编译错误或 sanitizer 报告；失败的是“加法可任意重排且逐位相同”的假设。修复不应只是固定一次偶然顺序，而应在第 16 章根据数据尺度选择更稳定算法，并用有来源的容差测试。

立即练习：把值的顺序改为 \texttt{\{1.0, 1.0e16, -1.0e16\}}。先预测，再验证正向结果为 0、反向结果为 1；元素集合没有变化。

\section{独立快照与输出任务}

在 WSL 中从项目根目录执行：

\begin{lstlisting}[language=bash]
cmake -S code/ch07 -B build/ch07
cmake --build build/ch07 -j
./build/ch07/ch07_ranges_pipeline code/ch07/data/trades.csv AAPL
\end{lstlisting}

输出必须与章首单行逐字一致。根项目的 CTest 还验证非法方向、缺失代码、惰性 view、浮点顺序、ASan 故障实验、核心练习与干净快照。Windows 可用同一 CMake 源目录选择 Visual Studio 生成器，运行时按所选配置进入 \texttt{Debug} 或 \texttt{Release} 子目录；正文的生命周期和算法规则不变。

\section{自测、编码任务与面试检查点}

\begin{enumerate}
  \item 为什么时间顺序成交默认适合 vector，而按代码仓位适合关联容器？
  \item \texttt{map::operator[]} 与 \texttt{at} 在键缺失时有何不同？
  \item \texttt{ranges::sort} 为什么作用于副本，而 \texttt{find\_if} 作用于原序列？
  \item view 的惰性从哪条输出得到证明？它拥有什么、借用什么？
  \item 返回借用局部容器的 view 为什么悬空？有哪些恢复方案？
  \item \texttt{back\_inserter} 在物化筛选结果时解决什么问题？
  \item 为什么 \texttt{accumulate} 的初值写 \texttt{0.0} 而不是 0？
  \item 相同浮点集合为何能归约出 1 和 0？本章应承诺什么、不承诺什么？
\end{enumerate}

编码任务：过滤固定成交中的 AAPL，按价格降序排序并只输出前两笔。可运行答案位于 \path{code/ch07/exercises/top_prices_solution.cpp}，精确输出必须为 \texttt{top=AAPL@102.00,AAPL@101.00}。

\begin{interviewcheck}
面试中若被问“ranges 是否一定比循环快”，先回答它主要表达组合与惰性，性能必须测量；再说明 view 的借用生命周期、物化时机、迭代器失效与浮点归约顺序。能指出这些边界，比背诵 adaptor 名称更可信。
\end{interviewcheck}

\section{本章小结}

容器选择来自访问、失效和所有权；ranges 算法直接处理范围。本章从左值容器建立的 view 惰性借用源数据，其他 view 的所有权取决于所包装的范围。需要排序、跨边界保存或固定所有权时应显式物化。归约接口不取消浮点舍入，数值稳定性留到第 16 章继续解决。
